\subsection{METRICS AND METRIC SPACES}

DEFINITION 1. A metric on a set X is a finite pseudometric d on X such that the relation \(d(x, y) = 0\) implies x = y. A metric space is a set X endowed with the structure defined by a given metric on X.

A metric space X is always considered as carrying the uniformity and the topology defined by the given metric on X.

Examples. 1) The Euclidean distance \(d(\pmb{x}, \pmb{y})\) (Chapter VI, § 2, no. 1) is a metric on real \(n\) -dimensional number space \(\mathbf{R}^n\) ; so are the functions

\[
\sup _ {1 \leqslant i \leqslant n} | x _ {i} - y _ {i} | \quad \text { and } \quad \sum_ {i = 1} ^ {n} | x _ {i} - y _ {i} |.
\]

All these metrics are equivalent (§ 1, no. 2).

\begin{enumerate}
\def\labelenumi{\arabic{enumi})}
\setcounter{enumi}{1}
\tightlist
\item
  On any set X the pseudometric \(d\) , defined by the relations \(d(x, x) = 0\) and \(d(x, y) = 1\) if \(x \neq y\) , is a metric. The uniformity it defines on X is the discrete uniformity.
\end{enumerate}

We have a definition equivalent to Definition 1 if we say that a metric is a finite pseudometric such that the uniformity defined by this pseudometric is Hausdorff; a finite pseudometric which is equivalent to a metric is therefore a metric.

Uniform spaces defined by a single pseudometric (which we may assume to be finite) can be reduced to metric spaces when the pseudometric is not a metric. Let \(f\) be such a pseudometric on a set \(\mathbf{X}\) , and let \(\mathcal{U}\) be the uniformity defined by \(f\) ; \(\mathcal{U}\) is not Hausdorff, and the intersection of the entourages of \(\mathcal{U}\) is the subset of \(\mathbf{X} \times \mathbf{X}\) defined by the equivalence relation \(f(x,y) = 0\) . Let \(\mathbf{R}\) denote this relation. If \(x \equiv x' \pmod{\mathbf{R}}\) , then by the triangle inequality we have

\[
f (x, y) \leqslant f (x, x ^ {\prime}) + f (x ^ {\prime}, y) = f (x ^ {\prime}, y)
\]

and similarly \(f(x', y) \leqslant f(x, y)\) , so that \(f(x, y) = f(x', y)\) ; in other words, \(f\) is a function compatible (in \(x\) and \(y\) ) with the equivalence relation \(\mathbf{R}\) (Set Theory, \(\mathbf{R}\) , § 5, no. 7). Let \(\overline{f}\) be the function induced by \(f\) on the quotient set; \(\overline{f}\) is defined on \((\mathbf{X}/\mathbf{R}) \times (\mathbf{X}/\mathbf{R})\) , and if \(x\) and \(y\) are any two points of \(\mathbf{X}\) and if \(\dot{x}\) and \(\dot{y}\) denote the equivalence classes (mod \(\mathbf{R}\) ) of \(x\) and \(y\) respectively, then we have \(\overline{f}(\dot{x}, \dot{y}) = f(x, y)\) . It follows immediately that \(\overline{f}\) is a metric on \(\mathbf{X}/\mathbf{R}\) ; it is called the metric associated with the pseudometric \(f\) ; furthermore, the uniformity it defines on \(\mathbf{X}/\mathbf{R}\) is precisely the Hausdorff uniformity associated with \(\mathcal{U}\) by the definition of this uniformity (Chapter II, § 3, no. 8, Remark). Thus, by passing to a suitable quotient space, the uniform structure defined by a single pseudometric can be reduced to the structure of a metric space.

Proposition 1 of § 1, no. 3 determines the structure of the completion of a metric space:

PROPOSITION 1. Let X be a metric space and let d be its metric. If \(\hat{X}\) is the completion of X (with respect to the uniformity defined by \(d\) ), the function \(d\) can be extended by continuity to \(\hat{X} \times \hat{X}\) ; the extended function \(\overline{d}\) is a metric on \(\hat{X}\) , and the uniformity of \(\hat{X}\) coincides with that defined by the metric \(\overline{d}\) .

Proposition 1 of § 1, no. 3 shows that \(\overline{d}\) is a finite pseudometric on \(\hat{\mathbf{X}}\) , and that the uniformity defined by \(\overline{d}\) on \(\hat{\mathbf{X}}\) is the uniformity obtained by completion; since this latter uniformity is Hausdorff, \(\overline{d}\) is a metric. Whenever we consider the completion of a metric space \(\mathbf{X}\) as a metric space, it is always to be understood that the metric on \(\hat{\mathbf{X}}\) is that obtained by extending the metric on \(\mathbf{X}\) by continuity.

